\section{Complete Verification}
\label{sec:comp-verification}
When incomplete methods yield an inconclusive result, complete verifiers are required to obtain a definitive decision. 
These verifiers typically employ a Branch-and-Bound (BaB) framework~\cite{bunelUnifiedViewPiecewise2018,wang2021beta}. 
Conceptually, BaB resolves the nonlinearities by systematically splitting the search space, 
either by branching over the input domain or by splitting unstable intermediate neurons. 
Within each branch of the search tree, an incomplete verifier evaluates the lower bound $\underline{\psi}_{\mathrm{sub}}$:
\begin{itemize}
    \item \textbf{Pruning:} 
        Subdomains that satisfy $\underline{\psi}_{\mathrm{sub}} > 0$ are formally certified and discarded from further search.
    \item \textbf{Branching / Falsification:} 
        Inconclusive subdomains with $\underline{\psi}_{\mathrm{sub}} \le 0$ are partitioned further, 
        or evaluated to find a concrete counterexample.
\end{itemize}
By iteratively adding branching constraints, the approximation error is systematically reduced across all subdomains. 
For piecewise-linear neural networks, this guarantees that complete verification terminates in a finite number of branching 
steps to yield a definitive decision~\cite{bunelUnifiedViewPiecewise2018,wang2021beta}.



\subsection{\texorpdfstring{$\alpha,\beta$}{α,β}-CROWN}
\label{subsec:alpha-beta-crown}
When applying Branch-and-Bound to neural networks with piecewise-linear activations such as ReLU, 
each branching step results in split constraints for unstable neurons. Specifically, an unstable neuron 
$j$ in layer $l$ is either forced into its active linear phase $z_j^{(l)} \ge 0$ or into its inactive phase 
$z_j^{(l)} < 0$~\cite{wang2021beta}. While classical LP solvers handle these constraints natively, 
they are computationally expensive and cannot be efficiently parallelized on GPUs~\cite{bunelUnifiedViewPiecewise2018}. 
However, unlike LP solvers, standard linear bound propagation frameworks such as CROWN cannot incorporate constraints 
on intermediate neurons, yielding loose bounds and redundant branching~\cite{xuFastCompleteEnabling2021,wang2021beta}.

To overcome this limitation, $\beta$-CROWN~\cite{wang2021beta} incorporates neuron split constraints directly into 
backward bound propagation using Lagrangian relaxation. By assigning non-negative Lagrange multipliers $\beta \ge 0$ 
to each split condition, constraint violations are penalized linearly during the backward pass.
The combination of the split constraints with the relaxation slopes $\alpha \in [0, 1]^{N_{\mathrm{unstable}}}$ from $\alpha$-CROWN 
yields the unified $\alpha,\beta$-CROWN framework. Propagating both relaxation parameters back to the input layer constructs 
the linear lower-bounding surrogate:
\begin{equation}
    w(\alpha, \beta)^\top x + c(\alpha, \beta) \;\le\; \psi(a^{(L)}(x)),
\end{equation}
where $w(\alpha, \beta) \in \mathbb{R}^{n_0}$ and $c(\alpha, \beta) \in \mathbb{R}$ denote the resulting surrogate coefficients. 
Minimizing this surrogate over the input domain $\mathcal{F}_{\mathrm{sub}} \subseteq \mathcal{F}$ of a given BaB subproblem 
defines the parameterized lower-bound objective:
\begin{equation}
    \mathcal{G}(\alpha, \beta) := \min_{x \in \mathcal{F}_{\mathrm{sub}}} \left( w(\alpha, \beta)^\top x + c(\alpha, \beta) \right).
    \label{eq:beta-g-def}
\end{equation}
For each BaB subdomain, the tightest lower bound is determined by jointly maximizing $\mathcal{G}(\alpha, \beta)$ 
over the relaxation slopes $\alpha$ and Lagrange multipliers $\beta$:
\begin{equation}
    \label{eq:beta-crown-opt}
    \underline{\psi}^\star_{\alpha,\beta} = \max_{\alpha \in [0, 1]^{N_{\mathrm{unstable}}}, \, \beta \ge 0} \mathcal{G}(\alpha, \beta) \;\le\; \psi^\star_{\mathrm{sub}}.
\end{equation}

Although $\mathcal{G}(\alpha, \beta)$ is concave w.r.t.\ $\beta$ for fixed linear relaxations, 
jointly optimizing over both the relaxation slopes $\alpha$ and the Lagrange multipliers $\beta$ leads to a non-concave optimization problem. 
Nonetheless, projected gradient-based optimizers can be effectively applied on GPUs to tighten the bounds. 
Because any feasible configuration $\left( \alpha \in [0, 1]^{N_{\mathrm{unstable}}}, \beta \ge 0 \right)$ corresponds to a valid dual relaxation, 
soundness is unconditionally preserved under early stopping, allowing fast and constraint-aware bounding in parallel.

Beyond branching over activation spaces, $\alpha,\beta$-CROWN can also branch directly over the input domain.
This input domain splitting is often more efficient for neural control systems with low-dimensional continuous state spaces. 
Furthermore, gradient-based adversarial attacks, such as PGD, may also be applied prior to the Branch-and-Bound search to identify potential 
counterexamples. If a counterexample is found, the Branch-and-Bound search is not needed and the verification terminates immediately. 

In summary, the verification methods presented in this chapter provide a flexible range of trade-offs between computational speed and bounding precision. 
CROWN offers efficient incomplete verification and can therefore be utilized in the Lyapunov co-training described in \cref{sec:lyapunov-learning}, 
where its lower bounds are included in the LiRPA-condition loss. In contrast, complete $\alpha,\beta$-CROWN verification provides formal certification 
by refining the verification domain. In the methodology presented in the following chapter, both approaches are combined: 
CROWN is used during training to promote a verifier-friendly Lyapunov formulation, 
while $\alpha,\beta$-CROWN is subsequently applied to formally certify the maximal region of attraction $\mathcal{V}_{\rho^\star}$.